JSON is een manier om data betekenis te geven. Het is onstaan op het web. Als we een webpagina hebben met informatie zoals:
\begin{lstlisting}[language=html]
<p>Jason Argonaut
Jupiterlaan 1
Heelverweg</p>
\end{lstlisting}
en we vragen die op, dan hebben we wel adres gegevens, maar we kunnen er niet zo veel mee.

De eerste oplossing was om daar XML van te maken:
\begin{lstlisting}[language=xml]
<voornaam>Jason</voornaam> <achternaam>Argonaut</achternaam>
<adres>Jupiterlaan 1</adres>
<plaats>Heelverweg</plaats>
\end{lstlisting}

Het zorgt voor meer betekenis van de ontvangen data, maar het wordt er niet leesbaarder op.

JSON heeft geprobeerd dit probleem op te lossen:
\begin{lstlisting}
{
    "Voornaam": "Jason",
    "Achternaam": "Argonaut",
    "Adres": "Jupiterlaan 1",
    "Plaats": "Heelverweg"
}
\end{lstlisting}


staat voor JavaScript Object Notation

Is voor gegevensuitwisseling
Menselijk leesbaar
Gemakkelijk te parsen door machines
Veel gebruikt in webapplicaties en API’s
